\documentclass[a4paper,12pt]{article}
\usepackage{iftex}
\ifPDFTeX
  \usepackage[T2A]{fontenc}
  \usepackage[utf8]{inputenc}
  \usepackage[scaled=0.97]{helvet}
  \renewcommand{\familydefault}{\sfdefault}
\else
  \usepackage{fontspec}
  \setmainfont{DejaVu Sans}
  \setsansfont{DejaVu Sans}
\fi
\usepackage[russian]{babel}
\usepackage{amsmath,amssymb,mathtools}
\usepackage{geometry}
\geometry{left=23mm,right=23mm,top=22mm,bottom=23mm}
\usepackage{microtype}
\usepackage{icomma}
\usepackage{setspace}
\setstretch{1.11}
\usepackage{enumitem}
\setlist{topsep=3pt,itemsep=3pt,parsep=0pt}
\usepackage{booktabs,tabularx,array}
\usepackage{xcolor}
\definecolor{accent}{HTML}{008080}
\definecolor{darkaccent}{HTML}{005F63}
\definecolor{ink}{HTML}{30383B}
\definecolor{muted}{HTML}{657177}
\usepackage{tikz}
\usetikzlibrary{arrows.meta,positioning,shapes.geometric,calc}
\usepackage{fancyhdr}
\usepackage[colorlinks=true,linkcolor=darkaccent]{hyperref}
\pagestyle{fancy}
\fancyhf{}
\fancyhead[L]{\footnotesize\color{muted} Показательные неравенства}
\fancyhead[R]{\footnotesize\color{muted} Николь Саркисьянц}
\fancyfoot[C]{\footnotesize\color{muted}\thepage}
\renewcommand{\headrulewidth}{0.25pt}
\renewcommand{\footrulewidth}{0pt}
\setlength{\headheight}{15pt}
\setlength{\parindent}{0pt}
\setlength{\parskip}{4pt}
\newcommand{\idea}[1]{\textcolor{darkaccent}{\textbf{#1}}}
\newcommand{\ans}[1]{\textcolor{darkaccent}{\textbf{Ответ:}}\ #1}
\newcommand{\sect}[1]{\section*{\color{darkaccent}#1}}
\ifPDFTeX\pdfmapfile{+cm-super-extended-t2a.map}\fi
\begin{document}
\thispagestyle{empty}
\begin{center}
  \vspace*{13mm}
  {\large\color{darkaccent} МАТЕМАТИКА\par}
  \vspace{10mm}
  {\fontsize{26}{30}\selectfont\color{ink} Чек-лист\par}
  \vspace{9mm}
  {\fontsize{22}{27}\selectfont\color{accent} Показательные неравенства\par}
  \vspace{6mm}
  {\Large\color{accent} Замена переменной и метод интервалов\par}
  \vspace{12mm}
  {\large\color{ink} Подготовка к ЕГЭ по профильной математике\par}
  \vfill
  {\large Николь Саркисьянц\par}
  \vspace{8mm}
  {\large 26 октября 2026 г.\par}
  \vspace{11mm}
  {\normalsize Лёвин Артём Александрович\par}
  {\normalsize эксклюзивно для Николь Саркисьянц\par}
  \vspace{10mm}
  \begin{tikzpicture}
    \draw[accent,line width=1pt] (0,0) .. controls (1.0,.65) and (2.0,-.65) .. (3,0)
       .. controls (4,.65) and (5,-.65) .. (6,0)
       .. controls (7,.65) and (8,-.65) .. (9,0);
    \fill[darkaccent] (4.5,0) circle (1.8pt);
  \end{tikzpicture}
  \vspace{2mm}
\end{center}
\clearpage

\sect{1. Что нужно уметь}
После повторения темы проверьте, что Вы можете выполнить каждый шаг самостоятельно.
\begin{itemize}[label=$\square$,leftmargin=9mm]
  \item Выделить повторяющееся выражение в показателе и ввести замену $t$.
  \item Представить степени через два основания и объединить одинаковые показатели.
  \item Разделить неравенство на заведомо положительную степень, сохранив знак.
  \item Решить квадратное неравенство с условием $k>0$, выполнить обратную замену.
  \item Решить полученные дробно-рациональные неравенства методом интервалов.
\end{itemize}

\sect{2. Правила, от которых зависит знак}
Показательная функция $a^u$ определена и \emph{строго положительна} при $a>0$. При $a\ne1$ её монотонность определяется основанием:
\[
\begin{array}{ll}
 a>1: & a^{f(x)}>a^{g(x)}\ \Longleftrightarrow\ f(x)>g(x),\\[4pt]
 0<a<1: & a^{f(x)}>a^{g(x)}\ \Longleftrightarrow\ f(x)<g(x).
\end{array}
\]
\idea{Запомните:} при сравнении степеней с основанием между нулём и единицей знак неравенства меняется.

Степени удобно преобразовывать по формулам
\[
 (ab)^t=a^t b^t,\qquad a^{2t}=(a^t)^2,\qquad
 \frac{a^t}{b^t}=\left(\frac ab\right)^t\quad(a,b>0).
\]
Например,
\[
 25^t=(5^t)^2,\qquad 15^t=3^t5^t,\qquad
 3^{2+2t}=9(3^t)^2.
\]

\sect{3. Допустимые преобразования}
\begin{itemize}[label=\textcolor{accent}{\textbullet},leftmargin=7mm]
\item \textbf{Деление.} Делить обе части неравенства на $(3^t)^2$ можно: это число положительно при любом действительном $t$.
\item \textbf{Произвольная переменная.} Деление на $x$ без выяснения его знака недопустимо: при отрицательном $x$ знак неравенства должен измениться. Если $x=0$, деление вообще невозможно.
\item \textbf{Новая переменная.} Если $k=(5/3)^t$, обязательно записать $k>0$.
\end{itemize}
\clearpage

\sect{4. Главный пример: две замены}
Решите неравенство
\begin{equation}\label{eq:main}
 3\cdot25^{1-\frac3x}-152\cdot15^{-\frac3x}
 +5\cdot3^{2-\frac6x}>0.
\end{equation}
\textbf{Шаг 1. ОДЗ и замена.} В показателях присутствует дробь, поэтому $x\ne0$. Обозначим
\[
 t=-\frac3x,\qquad x\ne0.
\]
Тогда получим
\[
 3\cdot25^{1+t}-152\cdot15^t+5\cdot3^{2+2t}>0.
\]
\textbf{Шаг 2. Приведение к основаниям $3$ и $5$.}
\begin{align*}
 3\cdot25\cdot25^t-152\cdot(3\cdot5)^t
    +5\cdot9\cdot3^{2t}&>0,\\
 75(5^t)^2-152\cdot3^t5^t+45(3^t)^2&>0.
\end{align*}
\textbf{Шаг 3. Деление на положительное выражение.}
Поскольку $(3^t)^2>0$, делим на него \emph{каждое} слагаемое и сохраняем знак:
\[
 75\left(\frac53\right)^{2t}
 -152\left(\frac53\right)^t+45>0.
\]
\textbf{Шаг 4. Вторая замена.} Пусть
\[
 k=\left(\frac53\right)^t>0.
\]
Неравенство принимает вид
\begin{equation}\label{eq:quad}
 75k^2-152k+45>0,\qquad k>0.
\end{equation}
\textbf{Шаг 5. Корни трёхчлена.}
\[
 D=(-152)^2-4\cdot75\cdot45=23104-13500=9604,
 \qquad\sqrt D=98.
\]
При устном поиске корня можно заметить $90^2<9604<100^2$, затем проверить $98^2=9604$.
\[
 k_1=\frac{152-98}{150}=\frac9{25}=\left(\frac35\right)^2,
 \qquad k_2=\frac{152+98}{150}=\frac53.
\]
\clearpage

\sect{5. Обратные подстановки и ответ}
\textbf{Шаг 6. Метод интервалов для $k$.} Поскольку старший коэффициент $75>0$, трёхчлен положителен снаружи своих корней. С учётом $k>0$:
\[
 k\in\left(0;\frac9{25}\right)\cup\left(\frac53;+\infty\right).
\]
\begin{center}
\begin{tikzpicture}[>=Latex,x=1cm]
 \draw[->,thick] (-.15,0)--(11,0) node[right] {$k$};
 \foreach \x/\txt in {0/0,3/{$\frac9{25}$},6/{$\frac53$}}{
  \draw (\x,0.11)--(\x,-0.11);
  \node[below=4pt] at (\x,0) {\txt};
 }
 \draw[accent,line width=2pt] (0.1,0)--(2.91,0);
 \draw[accent,line width=2pt,->] (6.11,0)--(10.1,0);
 \foreach \x in {0,3,6}{\draw[fill=white,thick] (\x,0) circle (2.6pt);}
 \node[above] at (1.5,.13) {$+$};\node[above] at (4.5,.13) {$-$};\node[above] at (8.3,.13) {$+$};
\end{tikzpicture}
\end{center}
\textbf{Шаг 7. Вернёмся к $t$.} Основание $5/3>1$, поэтому при сравнении показателей знак сохраняется:
\[
\left[\begin{gathered}
\left(\frac53\right)^t<\left(\frac53\right)^{-2},\\
\left(\frac53\right)^t>\left(\frac53\right)^1
\end{gathered}\right.
\quad\Longleftrightarrow\quad
\left[\begin{gathered}t<-2,\\t>1.\end{gathered}\right.
\]
\textbf{Шаг 8. Вернёмся к $x$.} Подставим $t=-3/x$:
\[
\left[\begin{gathered}-\dfrac3x<-2,\\[4pt]-\dfrac3x>1.\end{gathered}\right.
\]
Приведём левую часть каждого неравенства к одной дроби:
\[
\left[\begin{gathered}\dfrac{2x-3}{x}<0,\\[7pt]\dfrac{-x-3}{x}>0.\end{gathered}\right.
\]
\textbf{Шаг 9. Метод интервалов.}
Для первой дроби критические точки $0$ и $3/2$. Её знак отрицателен на $(0;3/2)$.
Для второй дроби критические точки $-3$ и $0$. Её знак положителен на $(-3;0)$.
В обоих случаях $0$ исключаем, поскольку $x$ находится в знаменателе.

\ans{$\displaystyle x\in(-3;0)\cup\left(0;\frac32\right).$}

\vspace{2mm}
\idea{Контроль:} сначала неравенство решалось относительно $k$. После обеих обратных замен ответ должен быть записан через $x$.
\clearpage

\sect{6. Второй приём: одинаковая структура степеней}
Решите неравенство, разобранное на занятии:
\[
 2\cdot12^x+9^x>3\cdot16^x.
\]
Выделим основания $3$ и $4$:
\[
 2\cdot3^x4^x+(3^x)^2>3(4^x)^2.
\]
Разделим на $(3^x)^2>0$ и обозначим $k=(4/3)^x>0$:
\[
 2k+1>3k^2
 \quad\Longleftrightarrow\quad 3k^2-2k-1<0
 \quad\Longleftrightarrow\quad (3k+1)(k-1)<0.
\]
Корни $-1/3$ и $1$. Для решения квадратного неравенства нужен промежуток $-1/3<k<1$; с учётом $k>0$ остаётся $0<k<1$.
Поскольку $4/3>1$, получаем
\[
 \left(\frac43\right)^x<1=\left(\frac43\right)^0
 \quad\Longleftrightarrow\quad x<0.
\]
\ans{$x\in(-\infty;0)$.}

\sect{7. Что делать, если дискриминант отрицателен?}
Рассмотрим квадратный трёхчлен $4k^2-5k+4$:
\[
 D=(-5)^2-4\cdot4\cdot4=-39<0.
\]
Корней нет, поэтому знак трёхчлена постоянен и совпадает со знаком старшего коэффициента. Здесь $4>0$, следовательно,
\[
 4k^2-5k+4>0\qquad\text{для всех }k\in\mathbb R.
\]
\idea{Важно:} отрицательный дискриминант допустим для неравенства. Необходимо выяснить знак всего выражения; отсутствие корней само по себе не означает отсутствие решений.

\clearpage
\sect{8. Финальный этап: рациональная дробь}
Решите неравенство
\[
 \frac{x+1}{3x+5}\ge4,\qquad x\ne-\frac53.
\]
Переносим всё влево, приводим к общему знаменателю:
\[
 \frac{x+1-4(3x+5)}{3x+5}\ge0
 \quad\Longleftrightarrow\quad
 \frac{-11x-19}{3x+5}\ge0.
\]
Критические точки: $-19/11$ (нуль числителя, включаем) и $-5/3$ (нуль знаменателя, исключаем). Проверка знака на трёх промежутках даёт
\[
 \displaystyle x\in\left[-\frac{19}{11};-\frac53\right).
\]
\clearpage

\begin{center}
{\LARGE\color{darkaccent} Алгоритм решения показательного неравенства}
\end{center}
\vspace{12mm}
\begin{center}
\begin{tikzpicture}[
 >=Latex,
 action/.style={draw=darkaccent,rounded corners=3pt,align=center,text width=11.6cm,
 minimum height=1.11cm,inner sep=5pt,font=\small},
 decision/.style={draw=accent,diamond,aspect=3.6,align=center,text width=5.6cm,
 inner sep=1.5pt,font=\small},
 smallaction/.style={draw=darkaccent,rounded corners=3pt,align=center,text width=5.3cm,
 minimum height=1.48cm,inner sep=4pt,font=\small},
 arrow/.style={->,thick,color=darkaccent,shorten >=2pt}
 ]
 \node[action] (a) {Проверить ОДЗ; выделить повторяющийся показатель и при необходимости ввести $t$};
 \node[action,below=10mm of a] (b) {Привести все степени к двум основаниям $a$ и $b$; собрать $a^t$, $b^t$ и их квадраты};
 \node[action,below=10mm of b] (c) {Разделить на $(b^t)^2>0$, затем обозначить $k=(a/b)^t>0$};
 \node[decision,below=12mm of c] (d) {Есть действительные корни\ квадратного трёхчлена?};
 \node[smallaction,below left=16mm and -7mm of d] (l) {Да: найти корни и определить нужные интервалы знакопостоянства};
 \node[smallaction,below right=16mm and -7mm of d] (r) {Нет: определить постоянный знак по старшему коэффициенту};
 \coordinate (joint) at ($(l.south)!0.5!(r.south)+(0,-9mm)$);
 \node[action,below=15mm of joint] (e) {Учесть $k>0$. Выполнить обратную замену $k\to t$ и сравнить показатели};
 \node[action,below=10mm of e] (f) {Вернуть $t\to x$. Дробные неравенства привести к одной дроби и решить методом интервалов};
 \node[action,below=10mm of f] (g) {Объединить подходящие промежутки; проверить выколотые точки и записать ответ};
 \draw[arrow] (a)--(b);
 \draw[arrow] (b)--(c);
 \draw[arrow] (c)--(d);
 \draw[arrow] (d.south west)--node[near start,left,font=\scriptsize]{да}(l.north);
 \draw[arrow] (d.south east)--node[near start,right,font=\scriptsize]{нет}(r.north);
 \draw[thick,color=darkaccent] (l.south)--(l.south |- joint);
 \draw[thick,color=darkaccent] (r.south)--(r.south |- joint);
 \draw[arrow] (joint)--(e);
 \draw[arrow] (e)--(f);
 \draw[arrow] (f)--(g);
\end{tikzpicture}
\end{center}
\clearpage

\sect{Тренировочный блок}
\textcolor{muted}{Путь А. 10 заданий по нарастающей сложности. Во всех заданиях решить неравенство.}
\begin{enumerate}[label=\textbf{\arabic*.},leftmargin=11mm,itemsep=9pt]
\item $3^{x+1}>81.$
\item $\displaystyle\left(\frac14\right)^{2x-1}\le16.$
\item $4^x-5\cdot2^x+4<0.$
\item $9^x-4\cdot3^x+3\ge0.$
\item $4\cdot9^x-5\cdot6^x+4\cdot4^x>0.$
\item $3\cdot4^x-5\cdot6^x+2\cdot9^x<0.$
\item $9\cdot16^x-25\cdot12^x+16\cdot9^x<0.$
\item $\displaystyle\frac{x+2}{2x+3}>2.$
\item $\displaystyle\frac{2x+1}{x-4}\ge3.$
\item $\displaystyle15\cdot25^{-\frac1x}-34\cdot15^{-\frac1x}
 +15\cdot9^{-\frac1x}>0.$
\end{enumerate}
\vspace{7mm}
\idea{Самопроверка перед ответом}
\begin{itemize}[label=\textcolor{accent}{\textbullet},leftmargin=7mm]
\item У новой переменной $k$ учтено условие $k>0$.
\item При основании меньше единицы знак сравнения показателей изменён.
\item Дроби приведены к общему знаменателю без неразрешённого умножения на $x$.
\item Нули знаменателей исключены; при нестрогом знаке нули числителя включены.
\item Ответ дан относительно $x$ как объединение нужных промежутков.
\end{itemize}
\clearpage

\sect{Ответы к тренировочному блоку}
\renewcommand{\arraystretch}{1.19}
\small
\begin{tabularx}{\linewidth}{@{}>{\bfseries}p{14mm}X@{}}
\toprule
№ & Ответ \\
\midrule
1 & $x\in(3;+\infty)$ \\
2 & $x\in\left[-\frac12;+\infty\right)$ \\
3 & $x\in(0;2)$ \\
4 & $x\in(-\infty;0]\cup[1;+\infty)$ \\
5 & $x\in\mathbb R$ \\
6 & $x\in(0;1)$ \\
7 & $x\in(0;2)$ \\
8 & $x\in\left(-\frac32;-\frac43\right)$ \\
9 & $x\in(4;13]$ \\
10 & $x\in(-1;0)\cup(0;1)$ \\
\bottomrule
\end{tabularx}
\vspace{8mm}
\sect{Короткая памятка}
\begin{tabularx}{\linewidth}{@{}p{53mm}X@{}}
\toprule
Сигнал в неравенстве & Ваше действие \\
\midrule
Повторяется дробь в показателе & Введите $t$ и запишите ОДЗ исходной переменной.\\
Два несводимых основания & Выделите степени $a^t$, $b^t$; разделите на положительную степень.\\
Получили квадратный трёхчлен & Введите $k>0$; найдите интервалы нужного знака.\\
$D<0$ & Определите постоянный знак по коэффициенту при $k^2$.\\
Вернулись к дроби с $x$ & Общий знаменатель, критические точки, метод интервалов.\\
Системы и совокупности & В системе берите пересечение; в совокупности --- объединение.\\
\bottomrule
\end{tabularx}
\vfill
\begin{center}
{\footnotesize\color{muted} Лёвин Артём Александрович эксклюзивно для Николь Саркисьянц}
\end{center}
\end{document}
